\section{Limitations and threats to validity}
\label{sec:limitations}

\subsection{Construct validity}

The most significant limitation is that the offline harness \emph{models}
pipeline latency rather than measuring it in a live cloud account. Our
stream-path figure of \SI{4}{\second} and batch-path figure of \SI{150}{\second}
are plausible for the services in question but they are parameters, not
observations. We mitigate this in three ways: the parameters are named rather
than buried; the sensitivity analysis reports results across a wide range of
them; and the deployed functions import the identical rule package, so a reader
with an account can substitute measured latencies without changing the detection
logic. Absolute latencies should nonetheless be read as \emph{lower bounds
conditioned on the modelled pipeline}, and the relative comparisons --- stream
against batch, threshold against threshold --- are the durable results.

The impact measure is coarse. Counting ``actions that succeeded'' treats an
executed thruster firing and a read object as one unit each, which is why we
report critical commands and exfiltrated volume separately. A severity-weighted
composite would be more informative and would also be more arbitrary; we chose
transparency.

\subsection{Internal validity}

Ground-truth labels are written by the attack scenarios themselves, so an error
in a scenario would produce a systematically wrong score. We guard against the
most dangerous form of this --- a rule reading the label it is being scored
against --- with a test that fails if the label appears anywhere in the rule
module. Attribution by evidence causality rather than by time window removes the
other common failure, in which unrelated benign alerts during an attack window
are counted as detections.

Attacks are separated by at least \SI{20}{\hour} so that alerts attribute
cleanly. Real intrusions overlap, and a campaign of simultaneous attacks would
stress deduplication and analyst triage in ways this design does not measure.

\subsection{External validity}

The benign workload is synthetic. Its noise sources are realistic in kind ---
typos, roaming operators, link fades, a weekly bulk job --- but their rates are
chosen. The false-positive figures should therefore be read as
\emph{order-of-magnitude} results conditioned on that workload, not as a
prediction for a specific mission. A team adopting this pack should expect to
re-tune thresholds against their own traffic; the sweep tables show how much
movement each threshold buys.

The configuration is one spacecraft, one ground station, one cloud account.
Multi-tenant ground-station-as-a-service, cross-mission analytics, and
constellation-scale operations each introduce isolation and correlation questions
this study does not address.

The spacecraft model is deliberately shallow: circular two-body propagation, no
perturbations, no manoeuvre dynamics, no link budget in decibels. This is
sufficient because the detections depend on pass timing, physical envelopes and
frame outcomes, all of which the model reproduces --- but it means the testbed
cannot be used to study attacks whose detection depends on finer physics, such as
subtle orbit-determination manipulation.

\subsection{What we did not evaluate}

Availability attacks (jamming, uplink flooding) are excluded because they are
detected by link metrics rather than security telemetry. Supply-chain compromise
of flight software, physical attacks on the ground station, and adversaries who
obtain the SDLS key are excluded by the threat model. Analyst experience --- the
question of whether \ResFPDay{} false positives per day is actually tolerable in
a given SOC --- is a human-factors question our method cannot answer.
